\section{Problem Setting and Data}\label{sec:setting}

\subsection{Setting}
An \emph{entity} is an account (Berka) or a (user, card) pair (TabFormer). The \emph{personal history size} $n$ is the number of the entity's own transactions available to fit a detector, taken from immediately before the evaluation window (\cref{sec:protocol}). A detector is unsupervised: it sees only the entity's past transactions, never labels.

We use \emph{behavioural anomaly} for a transaction that deviates from a pattern established by prior observations of the same entity. Behavioural anomalies, statistical outliers, suspicious transactions and confirmed fraud form no implication chain: a statistical outlier is a model artefact, a suspicious transaction is an institutional judgement involving policy and cost, and confirmed fraud requires evidence external to the transaction record. Most behavioural anomalies are benign (a holiday, a house move, a new job). A detector of behavioural anomalies identifies deviations from a baseline; it does not establish fraud or wrongdoing, and nothing in our labels is fraud.

We ask six questions. How does detection performance grow with $n$, and where does it saturate? Can a defensible sufficiency threshold be identified, and do different definitions of it agree? Does the requirement depend on the individual's behavioural variability? Is transaction count or temporal coverage the binding constraint? How do personal, population and hybrid models compare at each $n$? Do sophisticated detectors beat simple statistics under scarcity?

\subsection{Real Bank Data: PKDD'99 Berka}
The PKDD'99 Financial dataset \cite{berka1999pkdd} holds real, anonymised retail-banking transactions from a Czech bank: 1{,}056{,}320 transactions across 4{,}500 accounts from 1993-01-01 to 1998-12-31. Per-account depth has median 208, mean 234.7 and maximum 675; 3{,}883 accounts have at least 100 transactions, 1{,}835 at least 250 and 119 at least 500. Each transaction carries a date (day resolution), a credit/debit type, one of five operation modes, an amount, a balance, one of seven characterisation symbols (\texttt{k\_symbol}) and partner bank and account fields. The data have no fraud or anomaly labels.

Berka is the only real dataset we assessed that satisfies the requirements of per-entity identity, timestamps, sufficient per-entity depth and real behaviour. The most-cited card dataset (ULB) lacks a cardholder identifier, which makes per-user histories impossible; the IEEE-CIS competition data carry no entity key beyond a reverse-engineered competition artefact; and the remaining candidates (Sparkov, PaySim, BankSim, IBM AML) are synthetic with shallow per-entity histories or network-level tasks, while Bank Account Fraud contains application records rather than transaction sequences.

The primary host did not resolve from our execution environment, so files came from a GitHub mirror. The loader checks every one of the eight tables against the record counts published on the primary page and raises on mismatch; all eight matched. We additionally cross-checked every table byte-for-byte, after line-ending normalisation, against a second independent mirror; the SHA-256 digests are in the repository's data-integrity report. The data are released for an open research challenge and carry no machine-readable licence; we use them for research only and redistribute no raw records.

\subsection{Synthetic Card Data: TabFormer}
IBM's TabFormer credit-card dataset \cite{padhi2021tabular} is generator-produced, not real behaviour, so we never pool it with Berka or with the behavioural simulation. We use it for one narrow purpose: to test whether the curve shape found on one 1990s retail-banking dataset reappears on a schema with merchant, category, geographic and minute-resolution timing channels. It holds 24{,}386{,}900 transactions from 2{,}000 users, independently counted by a full chunked pass over the downloaded file. The entity is a (user, card) pair, because a person can hold several cards with materially different usage; 6{,}139 distinct pairs exist and 4{,}538 have at least 560 transactions. We use amount, chip use (swipe, chip or online), merchant name, merchant city and state, merchant category code (109 distinct values), the error field and time to the minute, plus the generator's own fraud label (0.122\% prevalence). Only the primary GitHub LFS endpoint was downloaded, so provenance is weaker than Berka's two-mirror check.

\subsection{Controlled Simulation}
The real data cannot supply merchant, category or geographic channels, nor \emph{controlled} variability. We therefore simulate 240 individuals across six archetypes of 40 each. Each individual is defined by a measurable parameter vector rather than a label: log-amount dispersion, merchant count and Zipf exponent, category count, location count and home-location weight, inter-arrival Gamma shape, hour-of-day von Mises concentration, and a drift rate that rotates merchant and category distributions over time. The archetypes are \emph{regular} (log-amount dispersion 0.30, 5 merchants, 1 location), \emph{moderate}, \emph{variable} (dispersion 1.10, 60 merchants, 8 locations), \emph{traveller} (15 locations, home weight 0.35), \emph{seasonal} (drift rate 0.004) and \emph{burst} (Gamma shape 0.35, producing inactivity punctuated by clusters). Anomalies fall into six classes (amount, merchant, category, geographic, temporal and multivariate), each generated from a deliberately shifted distribution independent of any detector; ground truth is the generating branch, which no detector observes. Simulation results are never pooled with real-data results.
